\SourceLeaf{285}%
\begin{center}\begin{tabular}{r|ccc}
Anfangspunkt&&&\\
&$-a$&$-b$&$+c$\\
8&&&\\
&$-a$&&\\
9&&&\\
&$+a$&&\\
12&&&\\
&$-a$&&\\
14&&&\\
&$+a$&&\\
16&&&\\
&$+a$&$-b$&\\
17&&&\\
&$-a$&$-b$&\\
18&&&\\
&$+a$&$-b$&\\
19&&&\\
&$+a$&$-b$&$+c$\\
20&&&\\
&$+a$&$+b$&$+c$\\
21&&&\\
&$-a$&$+b$&$+c$\\
22&&&\\
&$-a$&$-b$&$+c$\\
Anfangspunkt&&&\\
\end{tabular}\end{center}
Aus der vorigen Analyse sieht man, dass das letzte Zeichen, das in
dieser Tabelle jeder Figur entspricht, notwendig dasselbe sein muss
wie das erste. Das heißt, dass die Zeichen $+-$, die an den einzelnen
Schnittpunkten vorkommen, einander aufheben, wenn man diejenigen ausnimmt, die an den
einzelnen Punkten vorkommen, an denen sich die Figuren bilden.
Die Summe \textls{aller} Zeichen ist folglich dieselbe wie die der letzteren allein und
zugleich gleich der Summe der Zeichen im Charakter des
Anfangspunkts gegenüber den einzelnen Figuren $a$, $b$, $c$, d.\,h. allen,
wenn man die zuletzt gebildete Figur $d$ nicht mitzählt.
\par\Needspace{12\baselineskip}
Der Charakter auf der einen Seite der Figur $d$ gegenüber den einzelnen
Figuren werde durch $\alpha a+\beta b+\gamma c+\delta d$ aus[gedrückt],
und [man] setze $\alpha+\beta+\gamma+\delta=\sigma$. Der absolute
Charakter dieser Stelle hingegen sei $=\chi$. Dann ist die gesamte
Seitenbiegung
\[
= -\sigma\,.\,360^0 - \frac{F-\chi S}{S}\,.\,720^0,
\]
